\documentclass[aspectratio=169,10pt,xcolor=table]{beamer}

\usepackage[utf8]{inputenc}
\usepackage[T1]{fontenc}
\usepackage[brazilian]{babel}
\usepackage[sfdefault]{FiraSans}
\usepackage{booktabs}
\usepackage{array}
\usepackage{multicol}
\usepackage{fontawesome5}
\usepackage[most]{tcolorbox}
\usepackage{tikz}
\usetikzlibrary{arrows.meta,positioning,calc}

\newcolumntype{L}[1]{>{\raggedright\arraybackslash}p{#1}}

% Evita quebrar palavras com hífen
\hyphenpenalty=10000
\exhyphenpenalty=10000

% ================= CORES =================
\definecolor{azul}{HTML}{15324E}
\definecolor{azulclaro}{HTML}{E8EEF4}
\definecolor{solar}{HTML}{F2A007}
\definecolor{laranja}{HTML}{D2601A}
\definecolor{cinza}{HTML}{6B7680}
\definecolor{cinzaclaro}{HTML}{EEF0F2}
\definecolor{fundo}{HTML}{F7F5F0}
\definecolor{verde}{HTML}{2E7D4F}
\definecolor{vermelho}{HTML}{B3261E}

% ================= TEMA =================
\setbeamertemplate{navigation symbols}{}
\setbeamercolor{background canvas}{bg=white}
\setbeamercolor{normal text}{fg=azul}
\setbeamercolor{structure}{fg=laranja}
\setbeamercolor{alerted text}{fg=laranja}
\setbeamerfont{alerted text}{series=\bfseries}

\setbeamertemplate{itemize item}{\raisebox{0.15em}{\textcolor{solar}{\rule{0.42em}{0.42em}}}}
\setbeamertemplate{itemize subitem}{\textcolor{laranja}{\guillemotright}}
\setbeamertemplate{enumerate item}{%
  \tikz[baseline=(n.base)]\node[circle,fill=solar,text=azul,inner sep=0pt,minimum size=1.35em,
    font=\scriptsize\bfseries] (n) {\insertenumlabel};}
\setbeamertemplate{itemize/enumerate body begin}{\raggedright\setlength{\itemsep}{0.3em}}
\setbeamertemplate{itemize/enumerate subbody begin}{\setlength{\itemsep}{0.1em}}

% Título do slide: nome da parte em cima e título embaixo
\setbeamertemplate{frametitle}{%
  \vskip0.6em
  {\tiny\bfseries\color{laranja}\MakeUppercase{\insertsectionhead}}\par
  \vskip-0.1em
  {\Large\bfseries\color{azul}\insertframetitle}\par
  \vskip-0.15em
  \textcolor{solar}{\rule{2em}{2pt}}
  \vskip-0.2em}

% Rodapé com barra de progresso
\setbeamertemplate{footline}{%
  \begin{tikzpicture}[remember picture,overlay]
    \fill[cinzaclaro] (current page.south west) rectangle ($(current page.south east)+(0,0.09)$);
    \fill[solar] (current page.south west) rectangle
      ($(current page.south west)+({\paperwidth*(\insertframenumber/\inserttotalframenumber)},0.09)$);
  \end{tikzpicture}%
  \leavevmode\hbox to\paperwidth{%
    \hspace{1.6em}{\tiny\color{cinza}Instalação de Sistemas Fotovoltaicos \;\textbar\; Wenes Gomes Aquino}%
    \hfill{\tiny\bfseries\color{cinza}\insertframenumber\,/\,\inserttotalframenumber}\hspace{1.6em}}%
  \vskip0.55em}

\setbeamersize{text margin left=1.6em,text margin right=1.6em}
\linespread{1.08}

% ================= CARTÕES =================
% \begin{card}[opções]{cor}{título} ... \end{card}
\newtcolorbox{card}[3][]{enhanced,
  colback=#2!5!white, colframe=#2, coltitle=#2,
  colbacktitle=#2!5!white,
  boxrule=0pt, toprule=2.6pt, arc=3pt, outer arc=3pt,
  left=7pt, right=7pt, top=3pt, bottom=6pt,
  toptitle=5pt, bottomtitle=0pt, titlerule=0pt,
  fonttitle=\bfseries\small, fontupper=\small, halign title=flush left,
  title={#3}, before upper={\raggedright\setlength{\parskip}{0.2em}}, #1}

% Destaque (frase de efeito)
\newtcolorbox{destaque}{enhanced, colback=fundo, colframe=solar,
  boxrule=0pt, leftrule=3.5pt, arc=0pt, outer arc=0pt,
  left=10pt, right=10pt, top=3pt, bottom=3pt, before skip=4pt, after skip=0pt,
  fontupper=\small\itshape, halign=center}
\newcommand{\frase}[1]{\vspace{0.1em}\begin{destaque}#1\end{destaque}}

% Ícone colorido em círculo
\newcommand{\icone}[2]{\tikz[baseline=-0.6ex]\node[circle,fill=#1,text=white,
  minimum size=1.7em,inner sep=0pt,font=\footnotesize]{#2};}

% Tabelas com cabeçalho colorido e linhas alternadas
\newcommand{\cabecalho}{\rowcolor{azul}}
\newcommand{\tcab}[1]{\textcolor{white}{\bfseries #1}}

% ================= TIKZ =================
\tikzset{
  seta/.style={-{Stealth[length=2.4mm,width=2mm]},line width=0.9pt,cinza},
  pill/.style={rounded corners=3pt,font=\scriptsize\bfseries,align=center,
    minimum height=2.5em,inner sep=2pt},
}

% Barra das etapas: concluídas em azul, atual em amarelo, futuras em contorno
\newcommand{\barraetapas}[1]{%
\begin{center}
\begin{tikzpicture}
  \foreach \t [count=\i] in {Projeto,Planeja-\\mento,Estrutura,Módulos,Elétrica,Inversor,Comissio-\\namento,Entrega}{
    \ifnum\i<#1 \node[pill,fill=azul,text=white,text width=1.4cm] (n\i) at (\i*1.78,0) {\t};
    \else\ifnum\i=#1 \node[pill,fill=solar,text=azul,text width=1.4cm] (n\i) at (\i*1.78,0) {\t};
    \else \node[pill,draw=azul!40,line width=0.7pt,fill=white,text=azul!70,text width=1.4cm] (n\i) at (\i*1.78,0) {\t};
    \fi\fi
  }
  \foreach \i [evaluate=\i as \j using int(\i+1)] in {1,...,7}{\draw[seta] (n\i.east) -- (n\j.west);}
\end{tikzpicture}
\end{center}}

% Fluxo horizontal simples: \fluxo{largura}{A,B,C}
\newcommand{\fluxo}[2]{%
\begin{center}
\begin{tikzpicture}
  \foreach \t [count=\i] in {#2}{
    \node[pill,fill=azul,text=white,text width=#1] (f\i) at (\i*#1+\i*0.55cm,0) {\t};
    \xdef\ultimo{\i}
  }
  \foreach \i in {2,...,\ultimo}{\pgfmathtruncatemacro\a{\i-1}\draw[seta] (f\a.east) -- (f\i.west);}
\end{tikzpicture}
\end{center}}

% Abertura de cada parte
\newcommand{\abertura}[3]{%
\begin{frame}[plain]
  \begin{tikzpicture}[remember picture,overlay]
    \fill[azul] (current page.south west) rectangle (current page.north east);
    \node[anchor=east,text=white!8!azul,font=\fontsize{150}{150}\selectfont\bfseries]
      at ($(current page.east)+(-0.6,-0.2)$) {0#1};
    \fill[solar] ($(current page.west)+(1.2,0.95)$) rectangle ++(1.2,0.09);
    \node[anchor=west,text=solar,font=\normalsize\bfseries] at ($(current page.west)+(1.15,1.45)$) {PARTE #1};
    \node[anchor=north west,text=white,font=\fontsize{30}{34}\selectfont\bfseries,text width=10cm,align=left]
      (tt) at ($(current page.west)+(1.15,0.7)$) {#2};
    \node[anchor=north west,text=white!75!azul,font=\normalsize,text width=9cm,align=left]
      at ($(tt.south west)+(0,-0.2)$) {#3};
  \end{tikzpicture}
\end{frame}}

\title{Instalação de Sistemas Fotovoltaicos}
\author{Wenes Gomes Aquino}
\date{2026}

\begin{document}

% ================= CAPA =================
\begin{frame}[plain]
  \begin{tikzpicture}[remember picture,overlay]
    \fill[azul] (current page.south west) rectangle (current page.north east);
    % sol
    \begin{scope}[shift={($(current page.north east)+(-2.5,-2.2)$)}]
      \fill[solar] (0,0) circle (0.85);
      \foreach \a in {0,30,...,330}{\draw[solar,line width=2.4pt,line cap=round] (\a:1.12) -- (\a:1.45);}
    \end{scope}
    % arranjo de módulos
    \begin{scope}[shift={($(current page.south east)+(-6.4,0.9)$)},yslant=0.28]
      \foreach \x in {0,1.4,2.8,4.2}{\foreach \y in {0,1.0}{
        \fill[white!14!azul] (\x,\y) rectangle ++(1.3,0.9);
        \draw[solar!55,line width=0.35pt,step=0.325] (\x,\y) grid ++(1.3,0.9);
        \draw[white!45!azul,line width=0.8pt] (\x,\y) rectangle ++(1.3,0.9);}}
    \end{scope}
    \fill[solar] ($(current page.west)+(1.2,2.05)$) rectangle ++(1.2,0.09);
    \node[anchor=south west,text=solar,font=\normalsize\bfseries] at ($(current page.west)+(1.15,2.25)$) {RESUMO DOS CONTEÚDOS};
    \node[anchor=north west,text=white,font=\fontsize{30}{35}\selectfont\bfseries,text width=12cm,align=left]
      at ($(current page.west)+(1.15,1.85)$) {Instalação de\\Sistemas Fotovoltaicos};
    \node[anchor=north west,text=white!80!azul,font=\small,text width=10cm,align=left]
      at ($(current page.west)+(1.15,-1.0)$)
      {\textbf{\textcolor{white}{Wenes Gomes Aquino}}\\[0.25em]
       Prof. Brunno Henrique Brito, D. Eng.\\
       Pós-graduação \textbar\ IFTO \textbar\ 2026};
  \end{tikzpicture}
\end{frame}

% ================= ROTEIRO =================
\section{Roteiro}
\begin{frame}{Roteiro da apresentação}
  \begin{columns}[c]
    \begin{column}{0.44\textwidth}
      \begin{enumerate}
        \item Introdução e segurança
        \item Estruturas de fixação
        \item Instalação dos módulos
        \item Instalação elétrica
        \item Inversor e comissionamento
      \end{enumerate}
    \end{column}
    \begin{column}{0.5\textwidth}
      \begin{card}{azul}{\faLightbulb[regular]\enspace Ideia central da disciplina}
        Instalar um sistema fotovoltaico não é só fixar módulos e conectar cabos.
        É executar um sistema \textbf{elétrico e mecânico} feito para funcionar por \textbf{décadas}.
      \end{card}
    \end{column}
  \end{columns}
  \vspace{0.6em}
  \barraetapas{9}
\end{frame}

% ================= PARTE 1 =================
\section{Introdução e segurança}
\abertura{1}{Introdução e segurança}{O que é uma instalação bem executada e quais riscos existem na obra.}

\begin{frame}{O que é uma instalação bem executada?}
  Uma boa instalação atende a \alert{cinco requisitos}:
  \vspace{0.6em}
  \begin{center}
  \begin{tikzpicture}[req/.style={rounded corners=4pt,fill=fundo,text width=2.05cm,
      minimum height=3.3cm,align=center,font=\footnotesize,inner sep=4pt,anchor=north}]
    \foreach \ic/\t/\d [count=\i] in {
      \faHardHat/Segurança/Não coloca pessoas em risco,
      \faChartLine/Desempenho/Entrega a geração esperada,
      \faShield*/Confiabilidade/Funciona por anos sem falhas,
      \faClipboardCheck/Conformidade/Segue projeto{,} normas e distribuidora,
      \faTools/Manutenção/Permite inspeção e troca de peças}{
      \node[req] (r\i) at (\i*2.7,0) {\icone{azul}{\ic}\\[0.45em]{\bfseries\small\t}\\[0.3em]\d};
      \node[circle,fill=solar,text=azul,font=\scriptsize\bfseries,inner sep=0pt,minimum size=1.3em]
        at (r\i.north east) {\i};
    }
  \end{tikzpicture}
  \end{center}
  \frase{O sistema é elétrico e mecânico ao mesmo tempo, e precisa operar por décadas.}
\end{frame}

\begin{frame}{Do projeto à entrega da usina}
  \barraetapas{1}
  \vspace{0.2em}
  \begin{columns}[T]
    \begin{column}{0.48\textwidth}
      \begin{card}[equal height group=pp]{azul}{\faDraftingCompass\enspace 1. Projeto: o que o instalador precisa ler}
        \begin{itemize}
          \item layout, número e potência dos módulos
          \item strings, inversor e estrutura
          \item cabos, proteções e aterramento
          \item pontos de conexão
        \end{itemize}
      \end{card}
    \end{column}
    \begin{column}{0.48\textwidth}
      \begin{card}[equal height group=pp]{azul}{\faCalendarCheck[regular]\enspace 2. Planejamento: antes de ir para a obra}
        \begin{itemize}
          \item verificar o acesso e avaliar o telhado
          \item conferir materiais e ferramentas
          \item ver as condições do clima
          \item definir a equipe e identificar riscos
        \end{itemize}
      \end{card}
    \end{column}
  \end{columns}
\end{frame}

\begin{frame}{Logística, instalação e inspeção}
  \begin{columns}[T]
    \begin{column}{0.31\textwidth}
      \begin{card}[equal height group=lg]{azul}{\faTruck\enspace 3. Logística}
        Como levar \textbf{30 módulos de 30\,kg} ao telhado sem danificar os módulos e sem arriscar a equipe?
      \end{card}
    \end{column}
    \begin{column}{0.31\textwidth}
      \begin{card}[equal height group=lg]{azul}{\faWrench\enspace 4. Instalação}
        \begin{enumerate}\footnotesize
          \item estrutura
          \item módulos
          \item cabeamento
          \item inversor
          \item proteções e aterramento
        \end{enumerate}
      \end{card}
    \end{column}
    \begin{column}{0.31\textwidth}
      \begin{card}[equal height group=lg]{azul}{\faSearch\enspace 5. Inspeção}
        Antes de energizar, conferir: conexões, polaridade, torque, cabos, proteções, aterramento, identificação e configuração.
      \end{card}
    \end{column}
  \end{columns}
\end{frame}

\begin{frame}{Os 4 grandes riscos da obra}
  \begin{columns}[T]
    \begin{column}{0.48\textwidth}
      \begin{card}[equal height group=r1]{vermelho}{\faUserInjured\enspace 1. Queda de altura}
        Escadas, bordas, telhados frágeis ou molhados, circulação sobre telhas.
      \end{card}
      \begin{card}[equal height group=r2]{vermelho}{\faBolt\enspace 2. Choque elétrico}
        \textbf{Lado CA:} rede, proteção e inversor.\\
        \textbf{Lado CC:} módulos, strings, proteção e inversor.
      \end{card}
    \end{column}
    \begin{column}{0.48\textwidth}
      \begin{card}[equal height group=r1]{vermelho}{\faFire\enspace 3. Arco elétrico}
        Conexão CC mal feita causa aquecimento, carbonização e incêndio.
      \end{card}
      \begin{card}[equal height group=r2]{vermelho}{\faExclamationTriangle\enspace 4. Outros riscos}
        Ferramentas, carga pesada, cortes, calor, chuva, vento, animais e insetos.
      \end{card}
    \end{column}
  \end{columns}
  \frase{Antes de instalar, pare e pergunte: o que pode dar errado nesta obra?}
\end{frame}

\begin{frame}{ABNT NBR 17193 (fevereiro de 2025)}
  Segurança contra \alert{incêndio e choque elétrico} em sistemas FV em edificações.
  Protege moradores, equipes de manutenção e bombeiros.
  \vspace{0.5em}
  \begin{columns}[T]
    \begin{column}{0.48\textwidth}
      \begin{card}[equal height group=n1]{azul}{\faPowerOff\enspace Desligamento rápido (RSD)}
        Desenergizar os circuitos dos módulos em emergência, de forma manual ou automática.
      \end{card}
      \begin{card}[equal height group=n2]{azul}{\faTachometer*\enspace Limites de tensão}
        Tensão reduzida a níveis seguros em até \textbf{30\,s} após o desligamento rápido.
      \end{card}
    \end{column}
    \begin{column}{0.48\textwidth}
      \begin{card}[equal height group=n1]{azul}{\faFireExtinguisher\enspace Proteção contra arco (AFCI)}
        Detectar e interromper arcos no lado CC para evitar princípio de incêndio.
      \end{card}
      \begin{card}[equal height group=n2]{azul}{\faRoute\enspace Acesso e rotas}
        Disposição dos cabos, rotas de fuga e acesso livre para os bombeiros no telhado.
      \end{card}
    \end{column}
  \end{columns}
\end{frame}

\begin{frame}{Análise Preliminar de Riscos (APR)}
  \begin{center}\small
  \renewcommand{\arraystretch}{1.55}
  \rowcolors{2}{fundo}{white}
  \begin{tabular}{>{\bfseries}L{3.6cm} L{3cm} L{2.4cm} L{4cm}}
    \cabecalho \tcab{Etapa} & \tcab{Perigo} & \tcab{Consequência} & \tcab{Medida preventiva}\\
    Acesso ao telhado & Queda & Lesão grave & Proteção contra queda\\
    Transporte do módulo & Queda de objeto & Ferimentos & Içamento adequado\\
    Instalação dos módulos & Queda & Lesão & Sistema de proteção\\
    Conexão CC & Choque ou arco & Queimadura & Procedimento adequado\\
    Perfuração & Contato com fiação & Choque & Verificação prévia\\
    Trabalho externo & Chuva ou vento & Acidente & Suspender a atividade\\
  \end{tabular}
  \end{center}
\end{frame}

\begin{frame}{Regras de ouro do instalador}
  \begin{enumerate}
    \item Nunca comece uma instalação sem entender o projeto.
    \item Nunca subestime o risco de queda.
    \item Nunca considere o lado CC desligado só porque o inversor está desligado.
    \item Nunca improvise materiais, conectores ou procedimentos.
    \item Nunca energize sem fazer uma inspeção final.
  \end{enumerate}
  \frase{O bom instalador não é o mais rápido. É o que entrega uma instalação segura, correta e que continua funcionando muitos anos depois.}
\end{frame}

% ================= PARTE 2 =================
\section{Estruturas de fixação}
\abertura{2}{Estruturas de fixação}{O sistema fotovoltaico começa na estrutura.}

\begin{frame}{Por que a estrutura é tão importante?}
  \begin{columns}[c]
    \begin{column}{0.52\textwidth}
      A estrutura precisa garantir:
      \begin{enumerate}
        \item segurança dos módulos e das pessoas
        \item resistência mecânica
        \item estabilidade contra o vento
        \item durabilidade
        \item compatibilidade com os módulos
      \end{enumerate}
      \frase{O módulo não deve só ficar preso. Deve ficar seguro por décadas.}
    \end{column}
    \begin{column}{0.42\textwidth}
      \centering
      {\small\bfseries O caminho das cargas}\\[0.3em]
      \begin{tikzpicture}[node distance=0.2cm,
          c/.style={pill,text width=4.4cm,minimum height=1.55em,fill=azulclaro,text=azul}]
        \node[c,fill=solar] (a) {MÓDULO};
        \node[c,below=of a] (b) {GRAMPO};
        \node[c,below=of b] (c) {TRILHO};
        \node[c,below=of c] (d) {GANCHO / FIXADOR};
        \node[c,below=of d] (e) {TELHADO / ESTRUTURA};
        \node[c,below=of e,fill=azul,text=white] (f) {EDIFICAÇÃO};
        \foreach \x/\y in {a/b,b/c,c/d,d/e,e/f}{\draw[seta] (\x) -- (\y);}
      \end{tikzpicture}\\[0.3em]
      {\scriptsize\color{cinza} Toda força no módulo deve chegar com segurança à edificação.}
    \end{column}
  \end{columns}
\end{frame}

\begin{frame}{Principais tipos de estruturas}
  \begin{columns}[c]
    \begin{column}{0.5\textwidth}
      \small
      \renewcommand{\arraystretch}{1.6}
      \rowcolors{2}{fundo}{white}
      \begin{tabular}{>{\bfseries}L{2.6cm} L{3.6cm}}
        \cabecalho \tcab{Aplicação} & \tcab{Exemplo de uso}\\
        Telha cerâmica & Residências\\
        Telha metálica & Galpões\\
        Fibrocimento & Galpões e residências\\
        Laje & Edificações comerciais\\
        Solo & Usinas e sistemas grandes\\
      \end{tabular}
    \end{column}
    \begin{column}{0.44\textwidth}
      \begin{card}{verde}{\faCheckCircle\enspace Em todos os casos}
        \begin{itemize}
          \item fixar na estrutura do telhado, e não só na telha
          \item usar o fixador certo para cada cobertura
          \item vedar as perfurações
          \item seguir o manual do fabricante
        \end{itemize}
      \end{card}
    \end{column}
  \end{columns}
\end{frame}

\begin{frame}{Vento e distância entre ganchos}
  \begin{columns}[c]
    \begin{column}{0.33\textwidth}
      \begin{card}{azul}{\faWind\enspace Regra prática}
        O Brasil é dividido em \textbf{5 regiões de vento}.\\[0.3em]
        Quanto \textbf{mais vento}, \textbf{menor} a distância entre os ganchos.
      \end{card}
    \end{column}
    \begin{column}{0.62\textwidth}
      \centering\small
      \renewcommand{\arraystretch}{1.5}
      \rowcolors{2}{fundo}{white}
      \begin{tabular}{>{\bfseries}c c c c}
        \cabecalho \tcab{Região} & \tcab{Vento (km/h)} & \tcab{Módulo 60 cél.} & \tcab{Módulo 72 cél.}\\
        I   & até 108    & 2,00 m & 1,80 m\\
        II  & 108 a 126  & 1,70 m & 1,60 m\\
        III & 126 a 144  & 1,50 m & 1,40 m\\
        IV  & 144 a 162  & 1,35 m & 1,25 m\\
        V   & 162 a 180  & 1,20 m & 1,10 m\\
      \end{tabular}\\[0.4em]
      {\scriptsize\color{cinza} Exemplo de fabricante (perfil de alumínio de 45\,mm). Sempre seguir o manual do produto usado.}
    \end{column}
  \end{columns}
\end{frame}

\begin{frame}{Distância entre trilhos}
  \begin{columns}[c]
    \begin{column}{0.46\textwidth}
      \centering
      \begin{tikzpicture}[scale=0.95,cota/.style={{Stealth[length=1.6mm]}-{Stealth[length=1.6mm]},thin,cinza}]
        \fill[azulclaro] (0,0) rectangle (2.6,4.4);
        \draw[azul,line width=1.4pt] (0,0) rectangle (2.6,4.4);
        \foreach \y in {0.55,1.1,...,3.9}{\draw[azul!25] (0,\y) -- (2.6,\y);}
        \draw[azul!25] (1.3,0) -- (1.3,4.4);
        \fill[solar] (-0.5,0.92) rectangle (3.1,1.08);
        \fill[solar] (-0.5,3.32) rectangle (3.1,3.48);
        \draw[cota] (3.5,0) -- (3.5,1.0) node[midway,right,font=\scriptsize,text=azul] {1/5 a 1/4 de L};
        \draw[cota] (3.5,1.0) -- (3.5,3.4) node[midway,right,font=\scriptsize,text=azul] {1/2 a 3/5 de L};
        \draw[cota] (3.5,3.4) -- (3.5,4.4) node[midway,right,font=\scriptsize,text=azul] {1/5 a 1/4 de L};
        \draw[cota] (-0.9,0) -- (-0.9,4.4) node[midway,left,font=\small\bfseries,text=azul] {L};
        \node[font=\scriptsize\bfseries,text=laranja,anchor=east] at (-0.55,3.4) {trilho};
      \end{tikzpicture}
    \end{column}
    \begin{column}{0.5\textwidth}
      \begin{itemize}
        \item $L$ é o comprimento do módulo.
        \item Cada trilho fica a cerca de \textbf{1/5 a 1/4 de $L$} da borda.
        \item Entre os dois trilhos ficam \textbf{1/2 a 3/5 de $L$}.
      \end{itemize}
      \frase{O módulo fica apoiado nos pontos indicados pelo fabricante, sem esforço excessivo.}
    \end{column}
  \end{columns}
\end{frame}

\begin{frame}{Estrutura para solo}
  \begin{columns}[c]
    \begin{column}{0.55\textwidth}
      \small
      \renewcommand{\arraystretch}{1.5}
      \rowcolors{2}{fundo}{white}
      \begin{tabular}{L{3.2cm} L{3.4cm}}
        \cabecalho \tcab{Estrutura fixa} & \tcab{Rastreador (tracker)}\\
        Mais simples & Mais complexo\\
        Menor custo & Maior custo\\
        Menor manutenção & Maior manutenção\\
        Inclinação fixa & Inclinação variável\\
        Mais comum & Projetos específicos\\
        \textcolor{cinza}{--} & Maior produção potencial\\
      \end{tabular}
    \end{column}
    \begin{column}{0.4\textwidth}
      \begin{card}{azul}{\faMountain\enspace Pontos de atenção}
        \begin{itemize}
          \item tipo do solo e fundação
          \item ação do vento
          \item corrosão e drenagem
          \item nivelamento e alinhamento
        \end{itemize}
      \end{card}
    \end{column}
  \end{columns}
  \frase{No solo, a estrutura não termina no módulo: ela começa na fundação.}
\end{frame}

\begin{frame}{Aterramento da estrutura e erros comuns}
  \begin{columns}[T]
    \begin{column}{0.48\textwidth}
      \begin{card}[equal height group=at]{azul}{\faLink\enspace Aterramento de painéis e estruturas}
        \begin{itemize}
          \item molduras e trilhos \textbf{equipotencializados}
          \item cabo de aterramento entre módulos e trilho
          \item grampos com dentes que furam o anodizado e fazem contato elétrico
          \item estrutura ligada ao aterramento do sistema
        \end{itemize}
      \end{card}
    \end{column}
    \begin{column}{0.48\textwidth}
      \begin{card}[equal height group=at]{vermelho}{\faTimesCircle\enspace Erros de instalação}
        \begin{itemize}
          \item fixar só na telha
          \item parafuso inadequado
          \item não vedar furos
          \item trilhos desalinhados
          \item apertar demais
          \item improvisar peças
          \item ignorar o manual
        \end{itemize}
      \end{card}
    \end{column}
  \end{columns}
  \frase{Um bom módulo e um bom inversor não compensam uma estrutura mal instalada.}
\end{frame}

% ================= PARTE 3 =================
\section{Instalação dos módulos}
\abertura{3}{Instalação dos módulos}{Recebimento, manuseio, fixação, conexões e organização dos cabos.}

\begin{frame}{Recebimento e manuseio dos módulos}
  \begin{columns}[T]
    \begin{column}{0.31\textwidth}
      \begin{card}[equal height group=rm]{azul}{\faBoxOpen\enspace Ao receber}
        \begin{itemize}
          \item modelo, potência e quantidade
          \item número de série
          \item embalagem
          \item vidro, moldura e caixa de junção
          \item cabos e conectores
        \end{itemize}
      \end{card}
    \end{column}
    \begin{column}{0.31\textwidth}
      \begin{card}[equal height group=rm]{azul}{\faDolly\enspace No transporte}
        \begin{itemize}
          \item evitar impactos e torções
          \item não arrastar
          \item não apoiar objetos em cima
          \item considerar peso e tamanho
        \end{itemize}
      \end{card}
    \end{column}
    \begin{column}{0.31\textwidth}
      \begin{card}[equal height group=rm]{azul}{\faArrowUp\enspace Na subida ao telhado}
        \begin{itemize}
          \item transporte em dupla
          \item içamento ou guincho
          \item cuidado com o vento
          \item comunicação na equipe
        \end{itemize}
      \end{card}
    \end{column}
  \end{columns}
  \frase{Um módulo moderno pode ter mais de 3\,m² e funcionar como uma ``vela'' no vento.}
\end{frame}

\begin{frame}{Antes do primeiro módulo e sequência de instalação}
  \begin{columns}[c]
    \begin{column}{0.44\textwidth}
      \begin{card}{verde}{\faTasks\enspace Checklist antes de começar}
        \begin{itemize}
          \item estrutura concluída e trilhos alinhados
          \item layout e orientação conferidos
          \item sequência de instalação definida
          \item espaçamentos do projeto
          \item cabos e conectores acessíveis
        \end{itemize}
      \end{card}
    \end{column}
    \begin{column}{0.5\textwidth}
      \centering
      {\small\bfseries Sequência prática}\\[0.4em]
      \begin{tikzpicture}[node distance=0.1cm,
          c/.style={pill,text width=5.6cm,minimum height=1.45em,fill=azulclaro,text=azul,align=left,font=\scriptsize}]
        \node[c] (1) {\textbf{1.} Posicionar o módulo};
        \node[c,below=of 1] (2) {\textbf{2.} Ajustar o alinhamento};
        \node[c,below=of 2] (3) {\textbf{3.} Fixar};
        \node[c,below=of 3] (4) {\textbf{4.} Aplicar o torque especificado};
        \node[c,below=of 4] (5) {\textbf{5.} Fazer as conexões};
        \node[c,below=of 5] (6) {\textbf{6.} Organizar os cabos};
        \node[c,below=of 6,fill=solar] (7) {\textbf{7.} Conferir a instalação};
      \end{tikzpicture}
    \end{column}
  \end{columns}
\end{frame}

\begin{frame}{Danos visíveis e invisíveis}
  \begin{columns}[T]
    \begin{column}{0.31\textwidth}
      \begin{card}[equal height group=dn]{vermelho}{\faBan\enspace O que evitar}
        \begin{itemize}
          \item pisar no módulo
          \item apoiar peso no vidro
          \item arrastar na estrutura
          \item forçar cabos e conectores
          \item usar como apoio de ferramentas
        \end{itemize}
      \end{card}
    \end{column}
    \begin{column}{0.31\textwidth}
      \begin{card}[equal height group=dn]{azul}{\faEye\enspace Visíveis e invisíveis}
        \textbf{Visíveis:} vidro quebrado, moldura deformada, caixa de junção danificada.\\[0.4em]
        \textbf{Invisíveis:} microfissuras, danos às células, falhas internas.
      \end{card}
    \end{column}
    \begin{column}{0.31\textwidth}
      \begin{card}[equal height group=dn]{laranja}{\faSearchPlus\enspace Causas de microfissuras}
        \begin{itemize}
          \item impactos e quedas
          \item torção
          \item pressão localizada
          \item pisoteamento
          \item manuseio inadequado
        \end{itemize}
      \end{card}
    \end{column}
  \end{columns}
  \frase{Nem todo dano feito na instalação pode ser visto a olho nu.}
\end{frame}

\begin{frame}{Conexões entre módulos e crimpagem MC4}
  \begin{columns}[T]
    \begin{column}{0.46\textwidth}
      \begin{card}[equal height group=mc]{azul}{\faPlug\enspace Pontos de atenção}
        \begin{itemize}
          \item conferir a polaridade
          \item seguir a sequência da string
          \item conectores MC4 compatíveis
          \item encaixe completo
          \item \textbf{``pulo do sapo''}: ligação alternada que reduz o cabo de retorno da string
        \end{itemize}
      \end{card}
    \end{column}
    \begin{column}{0.5\textwidth}
      \begin{card}[equal height group=mc]{azul}{\faTools\enspace Crimpagem do conector MC4}
        \begin{enumerate}
          \item cortar o cabo
          \item decapar no comprimento correto
          \item inserir o condutor no terminal
          \item crimpar com o alicate específico
          \item inserir o terminal no conector
          \item conferir o travamento
          \item inspecionar a montagem
        \end{enumerate}
      \end{card}
    \end{column}
  \end{columns}
\end{frame}

\begin{frame}{Organização dos cabos}
  \begin{columns}[T]
    \begin{column}{0.48\textwidth}
      \begin{card}[equal height group=oc]{vermelho}{\faTimes\enspace Errado}
        \begin{itemize}
          \item cabos pendurados
          \item cabos encostando no telhado
          \item cabos sujeitos à abrasão
        \end{itemize}
      \end{card}
    \end{column}
    \begin{column}{0.48\textwidth}
      \begin{card}[equal height group=oc]{verde}{\faCheck\enspace Correto}
        \begin{itemize}
          \item cabos bem fixados
          \item string organizada
          \item conectores protegidos
        \end{itemize}
      \end{card}
    \end{column}
  \end{columns}
  \vspace{0.5em}
  {\small\bfseries\color{laranja} Conclusões da parte}
  \begin{itemize}\small
    \item O módulo deve chegar intacto à estrutura e ser posicionado conforme o projeto.
    \item Evitar impactos e esforços mecânicos durante a instalação.
    \item Uma inspeção final evita problemas futuros.
  \end{itemize}
\end{frame}

% ================= PARTE 4 =================
\section{Instalação elétrica}
\abertura{4}{Instalação elétrica}{Cabos, proteções CC e CA, qualidade das conexões e aterramento.}

\begin{frame}{Visão geral da instalação elétrica}
  \begin{center}
  \begin{tikzpicture}[c/.style={pill,text width=1.75cm,minimum height=2.7em}]
    \node[font=\scriptsize\bfseries,text=laranja,anchor=east,align=right] at (-1.15,0) {\faSolarPanel\\LADO CC};
    \foreach \t [count=\i] in {Módulos FV,Cabos CC,String / arranjo,Proteção CC}{
      \node[c,fill=solar,text=azul] (a\i) at (\i*2.45-2.45,0) {\t};}
    \node[c,fill=azul,text=white] (a5) at (4*2.45,0) {Inversor};
    \foreach \i [evaluate=\i as \j using int(\i+1)] in {1,...,4}{\draw[seta] (a\i) -- (a\j);}
    \node[font=\scriptsize\bfseries,text=cinza,anchor=east,align=right] at (-1.15,-1.8) {\faPlug\\LADO CA};
    \foreach \t [count=\i] in {Proteção CA,Quadro elétrico,Medidor,Rede da distribuidora}{
      \node[c,fill=cinzaclaro,text=azul] (b\i) at (\i*2.45-2.45,-1.8) {\t};}
    \draw[seta,rounded corners=4pt] (a5.south) -- ++(0,-0.4) -| (b1.north);
    \foreach \i [evaluate=\i as \j using int(\i+1)] in {1,...,3}{\draw[seta] (b\i) -- (b\j);}
  \end{tikzpicture}
  \end{center}
  \vspace{0.3em}
  \begin{card}{azul}{\faFile*[regular]\enspace Do projeto para a obra: o que interpretar}
    Diagrama unifilar, quantidade de strings, identificação dos circuitos, cabos e proteções especificados, aterramento e caminho dos cabos.
  \end{card}
\end{frame}

\begin{frame}{Cabos de corrente contínua (CC)}
  \begin{columns}[T]
    \begin{column}{0.48\textwidth}
      \begin{card}[equal height group=cb]{laranja}{\faThermometerHalf\enspace Cabo comum \emph{versus} cabo FV}
        \renewcommand{\arraystretch}{1.3}
        \begin{tabular}{@{}L{0.43\linewidth} L{0.5\linewidth}@{}}
          \textcolor{vermelho}{\faTimes\ \textbf{Cabo comum}} & \textcolor{verde}{\faCheck\ \textbf{Cabo para FV}}\\
          pouca resistência UV & resistente a UV\\
          isolação inadequada & dupla isolação\\
          envelhece rápido & feito para área externa\\
        \end{tabular}
      \end{card}
    \end{column}
    \begin{column}{0.48\textwidth}
      \begin{card}[equal height group=cb]{verde}{\faCheckCircle\enspace Boas conexões}
        \begin{multicols}{2}
        \begin{itemize}
          \item conectores compatíveis
          \item crimpagem correta
          \item encaixe sem folgas
          \item proteção contra água
          \item fora do contato com o telhado
          \item polaridade conferida
        \end{itemize}
        \end{multicols}
      \end{card}
    \end{column}
  \end{columns}
  \vspace{0.6em}
  \begin{center}
  \begin{tikzpicture}[c/.style={pill,fill=vermelho!10,text=vermelho!80!black,text width=1.9cm}]
    \foreach \t [count=\i] in {Conector mal crimpado,Resistência de contato,Aquecimento,Arco elétrico}{
      \node[c] (e\i) at (\i*2.75,0) {\t};}
    \node[c,fill=vermelho,text=white] (e5) at (5*2.75,0) {\faFire\ Risco de incêndio};
    \foreach \i [evaluate=\i as \j using int(\i+1)] in {1,...,4}{\draw[seta,vermelho!60] (e\i) -- (e\j);}
  \end{tikzpicture}
  \end{center}
\end{frame}

\begin{frame}{Proteções CC e CA}
  \begin{columns}[T]
    \begin{column}{0.48\textwidth}
      \begin{card}[equal height group=pr]{solar}{\faSolarPanel\enspace Lado CC}
        \begin{itemize}
          \item seccionamento CC
          \item fusíveis ou disjuntor CC, quando aplicáveis
          \item DPS CC
          \item RSD (desligamento rápido)
        \end{itemize}
      \end{card}
    \end{column}
    \begin{column}{0.48\textwidth}
      \begin{card}[equal height group=pr]{azul}{\faPlug\enspace Lado CA}
        \begin{itemize}
          \item \textbf{DPS CA}: protege contra surtos de tensão
          \item \textbf{Disjuntor}: protege contra sobrecarga e curto-circuito
        \end{itemize}
      \end{card}
    \end{column}
  \end{columns}
  \frase{Todo sistema precisa de string box separada? Não. Muitos inversores já têm chave CC e DPS internos. Quem define é o projeto.}
\end{frame}

\begin{frame}{Separar circuitos CC e CA}
  \begin{columns}[c]
    \begin{column}{0.42\textwidth}
      Evitar o \alert{mesmo eletroduto} para:\\[0.6em]
      \icone{solar}{\faSolarPanel}\enspace \textbf{CC}: módulos $\rightarrow$ inversor\\[0.6em]
      \icone{azul}{\faPlug}\enspace \textbf{CA}: inversor $\rightarrow$ quadro e rede
    \end{column}
    \begin{column}{0.54\textwidth}
      \begin{card}{azul}{\faQuestionCircle\enspace Por quê?}
        \begin{itemize}
          \item facilita identificação e manutenção
          \item reduz o risco de erro de conexão
          \item não mistura circuitos de características diferentes
          \item facilita o seccionamento nas intervenções
          \item mantém a instalação organizada e conforme o projeto
        \end{itemize}
      \end{card}
    \end{column}
  \end{columns}
\end{frame}

\begin{frame}{Qualidade nas conexões elétricas}
  \begin{center}\small
  \renewcommand{\arraystretch}{1.5}
  \rowcolors{2}{fundo}{white}
  \begin{tabular}{>{\bfseries}L{5.2cm} L{7.2cm}}
    \cabecalho \tcab{Recurso} & \tcab{Função principal}\\
    Terminal tubular (ilhós) e olhal & Melhorar a conexão de cabos flexíveis\\
    Barramento (fase, neutro, terra) & Distribuir e conectar condutores\\
    Conector de compressão & União mecânica e elétrica robusta\\
    Terminal bimetálico & Transição entre cobre e alumínio\\
    Arruela dentada & Auxiliar o contato e a equipotencialização\\
    Composto antioxidante & Proteção em aplicações específicas\\
    Torquímetro & Garantir o torque especificado\\
  \end{tabular}
  \end{center}
\end{frame}

\begin{frame}{Bem executada \emph{versus} erros comuns}
  \begin{columns}[T]
    \begin{column}{0.48\textwidth}
      \begin{card}[equal height group=be]{verde}{\faCheckCircle\enspace Instalação bem executada}
        \begin{itemize}
          \item cabos organizados, curvas adequadas
          \item conectores protegidos
          \item circuitos separados e identificados
          \item eletrodutos bem instalados
          \item cabos longe de superfícies abrasivas
          \item sem emendas desnecessárias
          \item quadro organizado
        \end{itemize}
      \end{card}
    \end{column}
    \begin{column}{0.48\textwidth}
      \begin{card}[equal height group=be]{vermelho}{\faTimesCircle\enspace Erros mais comuns}
        \begin{itemize}
          \item cabo CC inadequado
          \item conector mal crimpado ou incompatível
          \item cabo exposto ao sol, ou apoiado no telhado
          \item DPS ausente ou mal instalado
          \item aterramento inadequado
          \item disjuntor ou cabo CA mal dimensionado
          \item quadro sem identificação
        \end{itemize}
      \end{card}
    \end{column}
  \end{columns}
\end{frame}

\begin{frame}{Checklist elétrico final (antes de energizar)}
  \begin{columns}[T]
    \begin{column}{0.48\textwidth}
      \begin{card}[equal height group=ck]{solar}{\faSolarPanel\enspace Lado CC}
        \begin{itemize}
          \item polaridade conferida
          \item tensão das strings conferida
          \item conectores inspecionados
          \item cabos íntegros
          \item seccionamento adequado
          \item DPS instalado corretamente
          \item aterramento e equipotencialização
        \end{itemize}
      \end{card}
    \end{column}
    \begin{column}{0.48\textwidth}
      \begin{card}[equal height group=ck]{azul}{\faPlug\enspace Lado CA}
        \begin{itemize}
          \item tensão conferida
          \item sequência de fases, quando aplicável
          \item disjuntor adequado
          \item cabos e DPS CA conferidos
          \item PE (terra) conectado
          \item conexões apertadas
          \item circuitos identificados
        \end{itemize}
      \end{card}
    \end{column}
  \end{columns}
\end{frame}

\begin{frame}{Medição do aterramento com terrômetro}
  \begin{columns}[T]
    \begin{column}{0.5\textwidth}
      \begin{card}[equal height group=tr]{azul}{\faRulerHorizontal\enspace Resistência de aterramento}
        {\footnotesize\color{cinza} Método da queda de potencial (NBR 15749)}
        \begin{itemize}
          \item $X$ = maior dimensão do aterramento
          \item haste de corrente (H) a $3X$
          \item haste de tensão (S) a $Y = 0{,}62 \cdot 3X$
          \item repetir com $Y$ deslocado $\pm 5\,\%$
          \item diferença menor que 10\,\% indica leitura válida
          \item antes, medir a tensão espúria
        \end{itemize}
      \end{card}
    \end{column}
    \begin{column}{0.46\textwidth}
      \begin{card}[equal height group=tr]{azul}{\faGlobeAmericas\enspace Resistividade do solo}
        {\footnotesize\color{cinza} Método de Wenner (NBR 7117-1)}
        \begin{itemize}
          \item 4 hastes em linha, com espaçamento $D$ igual
          \item $D$ = 1, 2, 4, 8, 16, 32 e 64\,m
        \end{itemize}
        \begin{center}
          \tcbox[colback=white,colframe=solar,boxrule=1pt,arc=3pt,left=8pt,right=8pt]{\large$\rho = 2\pi R D$}
        \end{center}
        {\scriptsize $\rho$: resistividade ($\Omega\cdot$m) \quad $R$: leitura ($\Omega$) \quad $D$: distância (m)}
      \end{card}
    \end{column}
  \end{columns}
\end{frame}

% ================= PARTE 5 =================
\section{Inversor e comissionamento}
\abertura{5}{Inversor e comissionamento}{Instalar, conectar, testar e entregar o sistema.}

\begin{frame}{Antes de instalar e onde instalar o inversor}
  \begin{columns}[T]
    \begin{column}{0.48\textwidth}
      \begin{card}[equal height group=iv]{azul}{\faClipboardList\enspace Conferir antes}
        \begin{itemize}
          \item modelo, potência e configuração do projeto
          \item quantidade de MPPTs
          \item faixa de tensão CC e corrente por MPPT
          \item tensão da rede e número de fases
          \item materiais, conectores e manual
          \item local previsto no projeto
        \end{itemize}
      \end{card}
    \end{column}
    \begin{column}{0.48\textwidth}
      \begin{card}[equal height group=iv]{verde}{\faMapMarker*\enspace Local adequado}
        \begin{itemize}
          \item protegido de chuva e intempéries
          \item ventilação e dissipação de calor
          \item espaço livre ao redor
          \item acesso fácil para manutenção
          \item perto das conexões, sem obstáculos
          \item suporte mecanicamente adequado
        \end{itemize}
      \end{card}
    \end{column}
  \end{columns}
  \frase{Um inversor em local quente e sem ventilação pode funcionar hoje e falhar amanhã.}
\end{frame}

\begin{frame}{Fixação e conexões do inversor}
  \fluxo{1.75cm}{Marcação,Furação,Fixar o suporte,Nivelamento,Instalação}
  \vspace{-0.6em}
  \begin{columns}[T]
    \begin{column}{0.44\textwidth}
      \begin{card}[equal height group=fx]{azul}{\faExclamationCircle\enspace Atenção na fixação}
        \begin{itemize}
          \item parede compatível com o peso
          \item parafusos e buchas adequados
          \item não obstruir entradas e saídas de ar
          \item respeitar os afastamentos do fabricante
        \end{itemize}
      \end{card}
    \end{column}
    \begin{column}{0.52\textwidth}
      \begin{card}[equal height group=fx]{solar}{\faExchange*\enspace Entradas e saídas}
        \begin{minipage}[t]{0.47\linewidth}
          \textbf{Entradas}
          \begin{itemize}
            \item strings CC
            \item PE (terra)
            \item comunicação
            \item bateria*
          \end{itemize}
        \end{minipage}\hfill
        \begin{minipage}[t]{0.5\linewidth}
          \textbf{Saídas}
          \begin{itemize}
            \item CA
            \item PE (terra)
            \item monitoramento
            \item cargas de emergência*
          \end{itemize}
        \end{minipage}\par
        {\scriptsize\color{cinza} * Em inversores híbridos ou off-grid.}
      \end{card}
    \end{column}
  \end{columns}
\end{frame}

\begin{frame}{Comissionamento}
  \alert{Comissionamento} é a verificação e validação da instalação \textbf{antes da entrega}.
  \vspace{0.5em}
  \begin{columns}[T]
    \begin{column}{0.235\textwidth}
      \begin{card}[equal height group=cm]{vermelho}{\faShield*\enspace Segurança}
        conexões, polaridade, aterramento e proteções
      \end{card}
    \end{column}
    \begin{column}{0.235\textwidth}
      \begin{card}[equal height group=cm]{laranja}{\faCubes\enspace Integridade}
        cabos, conectores, fixações e identificação
      \end{card}
    \end{column}
    \begin{column}{0.235\textwidth}
      \begin{card}[equal height group=cm]{verde}{\faCogs\enspace Funciona\-mento}
        inicialização, parâmetros, geração e comunicação
      \end{card}
    \end{column}
    \begin{column}{0.235\textwidth}
      \begin{card}[equal height group=cm]{azul}{\faClipboardCheck\enspace Conformidade}
        projeto, fabricante e requisitos da instalação
      \end{card}
    \end{column}
  \end{columns}
  \vspace{0.4em}
  \begin{center}
  \begin{tikzpicture}[c/.style={pill,text width=3cm}]
    \node[c,fill=cinzaclaro,text=azul] (a) at (0,0) {Instalação concluída};
    \node[c,fill=solar,text=azul] (b) at (4,0) {Comissionamento};
    \node[c,fill=azul,text=white] (c) at (8,0) {Entrega};
    \draw[seta] (a) -- (b); \draw[seta] (b) -- (c);
  \end{tikzpicture}
  \end{center}
\end{frame}

\begin{frame}{Checklist do comissionamento}
  \newcommand{\ok}{\textcolor{verde}{\faSquare[regular]}}
  \begin{columns}[T]
    \begin{column}{0.47\textwidth}
      \small
      \renewcommand{\arraystretch}{1.5}
      \rowcolors{2}{fundo}{white}
      \begin{tabular}{L{5cm} c}
        \cabecalho \tcab{Verificação} & \tcab{OK?}\\
        Inversor bem fixado & \ok\\
        Conexões CC conferidas & \ok\\
        Polaridade CC conferida & \ok\\
        Conexões CA conferidas & \ok\\
        Torque das conexões & \ok\\
        Cabos sem danos & \ok\\
      \end{tabular}
    \end{column}
    \begin{column}{0.47\textwidth}
      \small
      \renewcommand{\arraystretch}{1.5}
      \rowcolors{2}{fundo}{white}
      \begin{tabular}{L{5cm} c}
        \cabecalho \tcab{Verificação} & \tcab{OK?}\\
        Conectores bem encaixados & \ok\\
        Proteções instaladas & \ok\\
        Aterramento conferido & \ok\\
        Circuitos identificados & \ok\\
        Parâmetros conforme projeto & \ok\\
        Área liberada para energizar & \ok\\
      \end{tabular}
    \end{column}
  \end{columns}
\end{frame}

\begin{frame}{Energização e verificação após ligar}
  \fluxo{1.6cm}{Inspeção final,Energizar,Inicializar,Configurar,Conectar à rede,Verificar}
  \begin{columns}[T]
    \begin{column}{0.48\textwidth}
      \begin{card}[equal height group=en]{azul}{\faMicrochip\enspace No inversor}
        \begin{itemize}
          \item tensão e corrente CC
          \item potência instantânea
          \item tensão, corrente e frequência CA
          \item energia gerada, status e códigos de erro
        \end{itemize}
      \end{card}
    \end{column}
    \begin{column}{0.48\textwidth}
      \begin{card}[equal height group=en]{azul}{\faDesktop\enspace No monitoramento}
        \begin{itemize}
          \item inversor on-line e comunicando
          \item strings e MPPTs como esperado
          \item geração compatível com o momento
          \item sem alarmes
        \end{itemize}
      \end{card}
    \end{column}
  \end{columns}
  \frase{Se o inversor liga, mas um MPPT está sem geração, o comissionamento não está correto.}
\end{frame}

\begin{frame}{Erros na instalação do inversor}
  \begin{card}{vermelho}{\faTimesCircle\enspace Erros comuns}
    \begin{multicols}{2}
    \begin{itemize}
      \item sem ventilação adequada
      \item conectores CC mal encaixados
      \item polaridade invertida
      \item torque inadequado
      \item cabos tensionados nos conectores
      \item cabos sem identificação
      \item aterramento não conferido
      \item configuração incorreta
      \item monitoramento não configurado
      \item entrega sem checklist
    \end{itemize}
    \end{multicols}
  \end{card}
  \frase{O sistema não está pronto quando o inversor liga. Está pronto quando foi instalado, inspecionado, testado, configurado e documentado.}
\end{frame}

\begin{frame}{Resumo geral: do inversor à entrega}
  \begin{center}
  \begin{tikzpicture}[c/.style={pill,fill=fundo,text=azul,font=\small,text width=5.4cm,align=left,
      minimum height=2.2em,inner sep=5pt}]
    \foreach \t [count=\i] in {Conferir projeto e manual,Instalar e fixar o inversor,Conferir o aterramento,
        Executar as conexões CC,Executar as conexões CA,Conferir as proteções}{
      \node[c] at (0,-\i*0.92) {\textcolor{laranja}{\bfseries\i.}\enspace\t};}
    \foreach \t [count=\i,evaluate=\i as \k using int(\i+6)] in {Inspeção completa,Energização,Configuração,
        Testes e monitoramento,Registro e entrega}{
      \node[c] at (6.6,-\i*0.92) {\textcolor{laranja}{\bfseries\k.}\enspace\t};}
  \end{tikzpicture}
  \end{center}
\end{frame}

% ================= ENCERRAMENTO =================
\section{Conclusão}
\begin{frame}{Conclusão}
  \barraetapas{9}
  \vspace{0.2em}
  \begin{itemize}
    \item A instalação começa no \textbf{projeto} e no \textbf{planejamento}, com foco em segurança.
    \item A \textbf{estrutura} leva as cargas do módulo até a edificação.
    \item Os \textbf{módulos} exigem cuidado para evitar danos visíveis e invisíveis.
    \item A \textbf{instalação elétrica} depende de cabos certos, boas conexões e proteções.
    \item O \textbf{comissionamento} valida tudo antes da entrega.
  \end{itemize}
  \frase{Um sistema bem instalado é seguro, gera o esperado e funciona por décadas.}
\end{frame}

\begin{frame}[plain]
  \begin{tikzpicture}[remember picture,overlay]
    \fill[azul] (current page.south west) rectangle (current page.north east);
    \begin{scope}[shift={($(current page.center)+(0,1.5)$)}]
      \fill[solar] (0,0) circle (0.45);
      \foreach \a in {0,45,...,315}{\draw[solar,line width=2pt,line cap=round] (\a:0.62) -- (\a:0.85);}
    \end{scope}
    \node[text=white,font=\fontsize{34}{38}\selectfont\bfseries] at ($(current page.center)+(0,-0.1)$) {Obrigado!};
    \node[text=white!75!azul,font=\normalsize] at ($(current page.center)+(0,-1.0)$) {Wenes Gomes Aquino};
    \node[text=solar,font=\small] at ($(current page.center)+(0,-1.5)$) {Instalação de Sistemas Fotovoltaicos};
  \end{tikzpicture}
\end{frame}

\end{document}
